\bigskip

\textbf{Partie 1}

\medskip

Soit $g$ la fonction définie pour tout nombre réel $x$ de l'intervalle $]0~;~+ \infty[$ par:

\[g(x) = \dfrac{2 \ln x}{ x}.\]

\smallskip

\begin{enumerate}
\item On note $g'$ la dérivée de $g$. Démontrer que pour tout réel $x$ strictement positif : 

\[g'(x) = \dfrac{2 - 2\ln x}{x^2}.\]

\item On dispose de ce tableau de variations de la fonction $g$ sur l'intervalle $]0~;~+ \infty[$ :

\begin{center}
\psset{unit=1cm}
\begin{pspicture}(9,3)

\psframe(9,3)\psline(0,2.5)(9,2.5)
\psline(2.5,0)(2.5,3)
\uput[u](1.25,2.4){$x$}\uput[u](3,2.4){$0$}\uput[u](5.5,2.4){1}\uput[u](7,2.4){e}\uput[u](8.5,2.4){$+ \infty$}
\psline(2.95,0)(2.95,2.5)\psline(3.05,0)(3.05,2.5)
\uput[u](3.5,0){$- \infty$}\rput(5.5,1.25){0}\uput[d](7,2.5){$\frac{2}{\text{e}}$}\uput[u](8.8,0){0}
\rput(1.25,1.5){Variations}
\rput(1.25,1.){de $g$}
\psline(4,0.32)(5.3,1.1)\psline{->}(5.6,1.3)(6.8,2.)\psline{->}(7.2,2)(8.6,0.4)
\end{pspicture}
\end{center}

\smallskip

Justifier les informations suivantes lues dans ce tableau:

	\begin{enumerate}
		\item la valeur $\dfrac{2}{\text{e}}$ ;
		\item les variations de la fonction $g$ sur son ensemble de définition ;
		\item les limites de la fonction $g$ aux bornes de son ensemble de définition.
	\end{enumerate}
\item En déduire le tableau de signes de la fonction $g$ sur l'intervalle $]0~;~+ \infty[$.
\end{enumerate}

\bigskip

\textbf{Partie 2}

\medskip

Soit $f$ la fonction définie sur l'intervalle $]0~;~+ \infty[$ par

\[f(x) = [\ln (x)]^2.\]

 Dans cette partie, chaque étude est effectuée sur l'intervalle $]0~;~+ \infty[$.

\medskip

\begin{enumerate}
\item Démontrer que sur l'intervalle $]0~;~+ \infty[$, la fonction $f$ est une primitive de la fonction $g$.
\item À l'aide de la partie 1, étudier :
	\begin{enumerate}
		\item la convexité de la fonction $f$ ;
		\item les variations de la fonction $f$.
\end{enumerate}
\item 
	\begin{enumerate}
		\item Donner une équation de la tangente à la courbe représentative de la fonction $f$ au point d'abscisse e.
		\item En déduire que, pour tout réel $x$ dans $]0~;~\text{e}]$ :
		
\[[\ln (x)]^2 \geqslant \dfrac{2}{\text{e}}x - 1.\]

	\end{enumerate}
\end{enumerate}
\bigskip
